% 11. Implementacja
\section{Implementacja}\label{sec:implementacja}

\begin{instrukcja}[zalecane]
Technologie z wersjami, struktura repozytorium i model danych. Nie wklejajcie całego kodu - wystarczą
1-2 najciekawsze fragmenty z krótkim opisem (np. walidacja formularza, zmiana układu na wąskim ekranie).
\end{instrukcja}

\subsection{Technologie}
\begin{table}[H]\centering
\caption{Wykorzystane technologie}\label{tab:technologie}
\begin{tabularx}{\linewidth}{@{}p{3.6cm}p{3.2cm}L@{}}
\naglowektabeli \nt{Technologia} & \nt{Wersja} & \nt{Do czego} \\
Astro & 7.4 & statyczna strona, routing plików, układ \texttt{Layout.astro} \\
React + \texttt{@astrojs/react} & 19.3 / 7.1 & interaktywne wyspy: Start, lista, szczegóły \\
TypeScript & 6.0 & typy danych (\texttt{Termin}) i logika \\
motion & 14.12 & animacja licznika dni (wzór: Magic UI Number Ticker) \\
lucide-react & 1.72 & ikony typów terminów i przycisków \\
CSS (bez frameworka) & -- & tokeny kolorów i odstępów, układ szeroki i wąski \\
\bottomrule
\end{tabularx}
\end{table}

\subsection{Struktura repozytorium}
\noindent\begin{minipage}{\linewidth}
\begin{lstlisting}
MOIU-Planer-Sesji-I-Terminow/
|-- README.md
|-- src/
|   |-- pages/ (index, terminy, termin, dodaj = E1-E4)
|   |-- layouts/Layout.astro (ramy graficzne, menu, skip-link)
|   |-- components/ (Nav, Breadcrumbs, TerminCard, TypBadge,
|   |               ListaTerminow, NajblizszyTermin, SzczegolyTerminu,
|   |               NumberTicker, SpotlightCard)
|   |-- lib/ (store.ts, format.ts, validate.ts, url.ts)
|   |-- data/terminy.json (dane startowe)
|   `-- styles/global.css (tokeny CSS)
`-- docs/ (K1, dokumentacja)
\end{lstlisting}
\end{minipage}

\subsection{Model danych}
Dane startowe są w \texttt{src/data/terminy.json} (12 terminów). Przy pierwszym wejściu aplikacja kopiuje je do \texttt{localStorage} i dalej wszystkie operacje działają tam (\texttt{store.ts}). Przykładowy rekord:
\begin{lstlisting}[language=json]
{
  "id": "t3",
  "przedmiot": "Metody opracowywania interfejsow uzytkownika",
  "typ": "kolokwium",
  "data": "2026-12-09",
  "godzina": "12:00",
  "sala": "B-204",
  "zakres": "Wyklady 1-4",
  "notatki": "",
  "przypomnienie": "1d"
}
\end{lstlisting}

\subsection{Wybrany fragment kodu}
Na liście ten sam link zachowuje się różnie zależnie od szerokości ekranu (\texttt{TerminCard.tsx}):
\begin{lstlisting}
function klik(e) {
  if (onWybierz && matchMedia('(min-width: 48rem)').matches) {
    e.preventDefault();   // szeroki ekran: zostajemy na liscie
    onWybierz(t.id);      // i pokazujemy szczegoly w panelu obok
  }                       // waski ekran: zwykle przejscie do /termin?id=
}
\end{lstlisting}
Karta jest zwykłym odnośnikiem, więc działa bez JavaScriptu, z klawiatury i w nowej karcie przeglądarki.
Tylko na szerokim ekranie kliknięcie jest przechwytywane i zmienia panel szczegółów oraz adres strony
(\texttt{history.replaceState}).
